\begin{table}[htbp]
\centering
\caption{稳健性检验与敏感性诊断}
\label{tab:robustness}
\small
\resizebox{\textwidth}{!}{%
\begin{tabular}{lrrrrr}
\toprule
规格 & 系数 & 标准误 & $p$值 & 观测值 & 结论 \\
\midrule
按成员国生效时点 & 0.0144 & 0.0041 & <0.001 & 630,280 & 支持 \\
剔除2020—2021年 & 0.0151 & 0.0047 & 0.001 & 472,710 & 支持 \\
样本窗口2019—2024年 & 0.0107 & 0.0042 & 0.011 & 472,710 & 支持 \\
剔除美国伙伴 & 0.0098 & 0.0042 & 0.019 & 508,248 & 支持 \\
剔除离岸金融中心 & 0.0171 & 0.0042 & <0.001 & 482,840 & 支持 \\
伙伴国聚类 & 0.0136 & 0.0203 & 0.503 & 630,280 & 推断敏感性 \\
企业—伙伴国双向聚类 & 0.0136 & 0.0201 & 0.498 & 630,280 & 推断敏感性 \\
双重机器学习（交叉拟合） & 0.0077 & 0.0021 & 0.002 & 78,785 & 方法敏感性 \\
\bottomrule
\end{tabular}}
\begin{tablenotes}
\footnotesize 注：前五行是预先设定的样本或处理时点稳健性检验；“支持”表示系数与主规格同方向且在5\%水平显著。随后两行改变聚类层级，用于检验推断对政策暴露层级的敏感性。DML行的标准误为企业—伙伴关系对聚类的999次自助法标准误。
\end{tablenotes}
\end{table}
